\begin{frame}
    \frametitle{Подсчёт перестановок ножей $\sigma(i,j)$}
    Первый раскрой на $i$-й ПРС:
    \[
    \sigma(i,1) = |y_{i,1}| + 1
    \]
    Последующие раскрои ($j>1$):
    \[
    \sigma(i,j) = \sigma(i,j{-}1) + |y_{i,j}| - \max_n\!\Bigl\{0\le n \le \min\bigl\{|y_{i,j-1}|,\,|y_{i,j}|\bigr\} : \forall k{\le}n\;\; r^Y_{i,j-1,k}=r^Y_{i,j,k}\Bigr\}
    \]
    \vspace{0.2cm}
    \begin{itemize}
        \item Учитывается совпадение начальных форматов текущего и~предыдущего раскроев
        \item Совпадающие позиции не~требуют перестановки ножей
        \item Порядок форматов в~раскрое влияет на~значение $\tau(\overline{\overline{Y}})$
    \end{itemize}
\end{frame}
